% Version 3.1 December 2024
% Based on the official Springer Nature LaTeX authoring template.
% Draft status: internal R1 manuscript scaffold, not submission-ready.

\documentclass[pdflatex,sn-nature]{sn-jnl}

\usepackage{graphicx}
\usepackage{amsmath,amssymb,amsfonts}
\usepackage{booktabs}
\usepackage{array}

\raggedbottom

\begin{document}

\title[Interpretable variant perturbation diagnostics]{Sparse protein language model features for interpretable variant perturbation diagnostics}

\author*[1]{\fnm{Zipeng} \sur{Liu}}\email{liuzipeng@tju.edu.cn}
\author[1]{\fnm{Coauthors} \sur{to be added}}

\affil*[1]{\orgdiv{Affiliation to be updated}, \orgname{Institution to be updated}, \orgaddress{\city{City}, \country{Country}}}

\abstract{Protein language models and specialized variant-effect predictors provide strong scalar estimates of missense pathogenicity, but scalar scores rarely explain which molecular signal a variant perturbs. We evaluated sparse autoencoder features from ESM-2-3B as an interpretable perturbation representation for human variant analysis. Annotation-linked sparse features produced residue-level firing examples and improved the local ESM-2 masked-marginal log-likelihood baseline when combined with LLR on ClinVar and cancer holdout variants. However, AlphaMissense and gMVP remained stronger scalar pathogenicity predictors, a direct AlphaMissense+SAE ensemble gate failed, LOF/GOF/DN classification collapsed under protein-level holdout, and ProteinGym did not show a mean advantage over LLR. The current indel extension is limited to reconstructable protein-sequence edits. These results support sparse feature perturbations as a hypothesis-generation and residual-diagnostic framework, rather than as a finished general mechanism classifier or state-of-the-art pathogenicity predictor.}

\keywords{protein language models, sparse autoencoders, variant interpretation, missense variants, indels, mechanism diagnostics, ESM-2}

\maketitle

\section{Introduction}\label{sec:introduction}

Protein language models (PLMs) can score protein variants directly from sequence using masked-marginal or ensemble likelihood ratios \cite{lin2023esm,meier2021esm1v}. Specialized predictors such as AlphaMissense now provide strong proteome-wide estimates of human missense pathogenicity \cite{cheng2023alphamissense}. These scalar scores are useful for prioritization, but they do not usually explain whether a variant perturbs a catalytic site, structural core, binding interface, regulatory region or other functional axis.

Sparse autoencoders (SAEs) offer a way to inspect PLM activations as sparse features rather than dense hidden states. If an SAE feature can be associated with residue-level annotations, then the wild-type to mutant feature delta can be treated as an interpretable perturbation signature. The central question is not simply whether this representation improves pathogenicity AUC, but whether it provides information that is orthogonal to scalar scores and useful enough to guide variant review, residual diagnostics and future experimental design.

Here we report the R1 branch of BioCC, an ESM-2 sparse-feature workflow for variant perturbation diagnostics. We trained and analyzed SAEs over five ESM-2-3B layers, linked features to residue annotations, evaluated pathogenicity on ClinVar and cancer holdout cohorts, tested LOF/GOF/DN mechanism classification, benchmarked against ProteinGym and extended the scoring pipeline to reconstructable indels. The results are mixed and informative: sparse perturbations add signal to ESM-2 LLR and produce interpretable feature tables, but they do not outperform AlphaMissense or gMVP, do not improve AlphaMissense in direct ensemble tests, do not generalize as a robust cross-protein mechanism classifier and do not improve generic DMS fitness prediction over LLR.

\section{Results}\label{sec:results}

\subsection{Sparse features gain residue-level annotation coverage}

% Evidence: Research1/results/annotation_alignment/expanded_summary_firing_20260503.json; Research1/results/variant_effect/mechanism_feature_audit_firing_20260504.md
We trained ESM-2-3B SAEs on layers 19, 23, 27, 31 and 35 and matched feature firing positions to Swiss-Prot, Pfam, Gene Ontology and binding-site annotations. Saving residue-level firing positions increased the number of interpretable features across all production layers. At F1 $>$ 0.5, known-feature counts increased from 146 to 381 at layer 19, 135 to 301 at layer 23, 72 to 163 at layer 27, 49 to 86 at layer 31 and 32 to 45 at layer 35 (Table~\ref{tab:annotation}). The deepest layer improved but remains below the internal target of 60 known features.

\begin{table}[h]
\caption{Annotation alignment after saving firing positions. Known, partial and useful features use F1 thresholds of 0.5, 0.2 and 0.1, respectively.}\label{tab:annotation}
\begin{tabular}{@{}rrrrrr@{}}
\toprule
Layer & Known before & Known after & Partial after & Useful after & Firing-position features \\
\midrule
19 & 146 & 381 & 3483 & 7481 & 13868 \\
23 & 135 & 301 & 2992 & 6773 & 13832 \\
27 & 72  & 163 & 1725 & 4423 & 13456 \\
31 & 49  & 86  & 1135 & 3197 & 13649 \\
35 & 32  & 45  & 653  & 1935 & 12895 \\
\botrule
\end{tabular}
\end{table}

The feature audit produced 150 class-linked feature rows for DN, GOF and LOF labels across five layers. Each row records coefficient weight, feature kind, best annotation, Pfam or functional proxy and top residue-level firing examples. The table is suitable for manual residue triage and figure planning, but not for direct mechanistic claims without protein-specific biological review.

\subsection{Sparse perturbation features complement LLR but not AlphaMissense}

% Evidence: Research1/results/variant_effect/available_baseline_summary_20260507.md; Research1/results/variant_effect/t1a_final_available_baselines_no_primateai_20260510.md; Research1/results/variant_effect/grouped_pathogenicity_baselines_20260511.md
On 2,000 ClinVar missense variants, annotation-selected SAE features alone reached AUC 0.8782, similar to ESM-2 LLR at 0.8822. Combining SAE-LR and LLR by z-sum improved AUC to 0.9143. On the 101-variant cancer holdout, SAE-LR reached 0.9079, LLR reached 0.8978 and SAE+LLR reached 0.9193 (Table~\ref{tab:pathogenicity}). This supports complementarity between sparse perturbation signatures and masked-marginal likelihood, but not superiority over specialized missense predictors.

\begin{table}[h]
\caption{Pathogenicity performance on local R1 cohorts. Confidence intervals are bootstrap 95\% intervals.}\label{tab:pathogenicity}
\begin{tabular}{@{}llllr@{}}
\toprule
Cohort & Method & AUC & 95\% CI & n \\
\midrule
ClinVar2000 & SAE-LR & 0.8782 & [0.8619, 0.8923] & 2000 \\
ClinVar2000 & ESM-2 LLR & 0.8822 & [0.8677, 0.8960] & 2000 \\
ClinVar2000 & SAE+LLR & 0.9143 & [0.9010, 0.9259] & 2000 \\
ClinVar2000 & AlphaMissense & 0.9474 & [0.9377, 0.9567] & 1972 \\
ClinVar2000 & gMVP & 0.9369 & [0.9257, 0.9474] & 1788 \\
ClinVar2000 & ESM-1v & 0.9089 & [0.8952, 0.9212] & 1972 \\
CancerHoldout101 & SAE-LR & 0.9079 & [0.8471, 0.9590] & 101 \\
CancerHoldout101 & ESM-2 LLR & 0.8978 & [0.8273, 0.9561] & 101 \\
CancerHoldout101 & SAE+LLR & 0.9193 & [0.8617, 0.9646] & 101 \\
CancerHoldout101 & AlphaMissense & 0.9700 & [0.9244, 0.9988] & 101 \\
CancerHoldout101 & gMVP & 0.9400 & [0.8878, 0.9806] & 98 \\
CancerHoldout101 & ESM-1v & 0.8552 & [0.7768, 0.9211] & 101 \\
\botrule
\end{tabular}
\end{table}

AlphaMissense was stronger on both cohorts, with AUC 0.9474 on ClinVar2000 and 0.9700 on CancerHoldout101. gMVP was also stronger than SAE+LLR on both cohorts with available matches. A gene-bootstrap ClinVar2000 table preserved this ranking: AlphaMissense AUC 0.9474, gMVP 0.9369, ESM-1v 0.9089 and SAE-LR under gene-grouped cross-validation 0.7559. Thus R1 should not be framed as a state-of-the-art scalar pathogenicity predictor. The defensible novelty is interpretability, residual information and coverage of variant classes or mechanism questions that scalar missense predictors do not directly answer.

\subsection{Direct AlphaMissense ensembles fail the complementarity gate}

% Evidence: Research1/results/variant_effect/alphamissense_sae_ensemble_20260511.md
To test whether SAE features add scalar information beyond AlphaMissense, we ran a conservative grouped cross-validation ensemble on the 1,972 ClinVar2000 variants matched to AlphaMissense. AlphaMissense alone reached AUC 0.9474. SAE-LR under gene-symbol grouped cross-validation reached AUC 0.7559. Both tested ensembles were worse than AlphaMissense: AM+SAE stacking reached AUC 0.8542 and AM+SAE z-sum reached AUC 0.9210 (Table~\ref{tab:amensemble}).

\begin{table}[h]
\caption{AlphaMissense plus SAE ensemble gate on ClinVar2000 matched variants. Confidence intervals are bootstrap 95\% intervals; delta intervals use group bootstrap over gene symbols.}\label{tab:amensemble}
\begin{tabular}{@{}llll@{}}
\toprule
Method & AUC & 95\% CI & Delta vs AlphaMissense \\
\midrule
AlphaMissense & 0.9474 & [0.9373, 0.9568] & reference \\
SAE-LR group-CV & 0.7559 & [0.7341, 0.7771] & -0.1915 [-0.2171, -0.1660] \\
AM+SAE stack & 0.8542 & [0.8382, 0.8701] & -0.0932 [-0.1098, -0.0770] \\
AM+SAE z-sum & 0.9210 & [0.9083, 0.9319] & -0.0264 [-0.0344, -0.0189] \\
\botrule
\end{tabular}
\end{table}

This gate fails the scalar complementarity claim against AlphaMissense. SAE perturbations should therefore be used as an interpretation and residual-diagnostic layer rather than as an AlphaMissense ensemble improvement.

\subsection{Mechanism classification does not generalize across proteins}

% Evidence: Research1/results/variant_effect/mechanism_classifier_results_t0_protein_holdout_20260429.json; Research1/results/variant_effect/available_baseline_summary_20260507.md
We assembled 497 mechanism-labeled variants across 255 proteins, covering 80 DN, 120 GOF and 297 LOF labels. Under variant-level cross-validation, SAE-only features achieved macro-AUC 0.7471, whereas LLR alone was near random at 0.5054. Under protein-level holdout, SAE-only macro-AUC dropped to 0.5161 and SAE+LLR to 0.5164 (Table~\ref{tab:mechanism}).

\begin{table}[h]
\caption{LOF/GOF/DN mechanism classification. Protein-level holdout is the decision metric because variant-level splitting leaks protein-specific signatures.}\label{tab:mechanism}
\begin{tabular}{@{}lll@{}}
\toprule
Method & Variant-level macro-AUC & Protein-level macro-AUC \\
\midrule
SAE-LR & 0.7471 & 0.5161 \\
ESM-2 LLR & 0.5054 & 0.4642 \\
SAE+LLR & 0.7488 & 0.5164 \\
\botrule
\end{tabular}
\end{table}

The variant-level result indicates that sparse perturbations contain mechanism-like structure within related proteins. The protein-level result prevents a stronger claim: the current classifier does not reliably predict LOF, GOF and DN for unseen proteins.

\subsection{ProteinGym is a negative diagnostic for generic DMS fitness}

% Evidence: Research1/results/variant_effect/proteingym_benchmark_sae_signed_diagnostics_20260503.json
We evaluated the ProteinGym substitution benchmark, scoring 217 DMS assays and obtaining finite SAE values for 214 assays. ESM-2 LLR achieved mean Spearman rho 0.4341. Raw SAE disruption had mean rho -0.2314, consistent with disruption magnitude increasing as DMS fitness decreases. A sign-corrected SAE+LLR ensemble improved to 0.4047 but remained below LLR and beat LLR on only 33.6\% of usable assays (Table~\ref{tab:proteingym}).

\begin{table}[h]
\caption{ProteinGym diagnostic results across usable substitution assays.}\label{tab:proteingym}
\begin{tabular}{@{}llll@{}}
\toprule
Score & Assays & Mean Spearman rho & 95\% CI \\
\midrule
ESM-2 LLR & 214 & 0.4341 & [0.4091, 0.4594] \\
SAE disruption & 214 & -0.2314 & [-0.2588, -0.2035] \\
Negated SAE disruption & 214 & 0.2314 & [0.2034, 0.2592] \\
Legacy SAE+LLR & 214 & 0.1993 & [0.1737, 0.2244] \\
Sign-corrected SAE+LLR & 214 & 0.4047 & [0.3818, 0.4282] \\
\botrule
\end{tabular}
\end{table}

ProteinGym therefore bounds the method. SAE perturbation magnitude is not a generic DMS fitness predictor. It is better interpreted as a disruption or mechanism-like diagnostic whose utility should be tested in targeted biological contexts.

\subsection{Indel scoring extends the perturbation diagnostic beyond missense}

% Evidence: Research1/results/variant_effect/indel_records_supported_20260504_summary.json; Research1/results/variant_effect/indel_mechanism_predictions_20260504_summary.json; results/final_plan_readiness/final_plan_readiness_20260511.md
We staged 185,655 ClinVar indel records. The protein-HGVS reconstruction pipeline supported 18,897 records across all labels and 6,649 binary-label, length-compatible records for the existing scorer. The completed scored set contained 5,274 pathogenic and 1,375 benign labels across deletion, insertion, delins and duplication classes. The transferred mechanism classifier predicted 4,161 LOF, 1,354 GOF and 1,134 DN records. The SAE damage-score AUC for indel pathogenicity was 0.7735.

This result is preliminary because the affected-region SAE delta is a heuristic and remains below the target AUC of 0.85. It also does not yet support an approximately 80,000-record indel claim: the main unsupported reasons are missing UniProt sequence mappings and frameshifts that require transcript-aware reconstruction. It nevertheless supports a distinct use case: sparse perturbation signatures can be computed for reconstructable sequence edits that are not uniformly covered by missense-only scalar tools.

\section{Discussion}\label{sec:discussion}

R1 currently supports a conservative article centered on interpretable perturbation diagnostics. Sparse features add signal to ESM-2 LLR and produce residue-linked feature tables that can guide manual biological review. They do not yet provide a robust cross-protein mechanism classifier and should not be positioned as a replacement for AlphaMissense on scalar pathogenicity.

For the current no-wet-lab manuscript, the strongest defensible claim is a transparent diagnostic workflow: residue-linked SAE features, scalar benchmarks against accessible baselines, explicit negative gates and a bounded indel extension. Future assay work could test whether variants with similar AlphaMissense or LLR scores but divergent SAE perturbation signatures enrich experimentally separable mechanisms. That future direction should not be required for the current article's core evidence.

\section{Online Methods}\label{sec:methods}

\subsection{Sparse autoencoder training and annotation}

SAEs were trained on frozen ESM-2-3B activations from layers 19, 23, 27, 31 and 35. Feature annotations were generated by aligning top firing residue positions to Swiss-Prot, Pfam, Gene Ontology and BioLiP/PDB binding-site proxies. Known, partial and useful annotation tiers used F1 thresholds of 0.5, 0.2 and 0.1.

\subsection{Variant scoring}

For each missense variant, wild-type and mutant sequences were encoded by ESM-2 and the layer-specific SAEs. Feature deltas were concatenated across production layers. Logistic regression was trained on annotation-selected feature columns. ESM-2 LLR was computed as the masked-marginal log-likelihood ratio between mutant and wild-type residues. SAE+LLR used z-scored component scores.

\subsection{Cohorts and baselines}

The local pathogenicity cohort contained 2,000 balanced ClinVar missense variants. The cancer holdout contained 101 variants. AlphaMissense, gMVP and ESM-1v were matched by UniProt identifier and protein substitution when available. PrimateAI-3D was excluded from this analysis round because the gated dataset was unavailable. The AlphaMissense ensemble gate used gene-symbol grouped stratified cross-validation to reduce protein-family leakage.

\subsection{Mechanism labels}

LOF, GOF and DN labels were assembled from curated variant and gene-level mechanism resources. Multinomial logistic regression models were evaluated under both variant-level cross-validation and protein-level holdout. Protein-level holdout was treated as the decision metric.

\subsection{ProteinGym and indels}

ProteinGym substitution assays were evaluated by per-assay Spearman correlation between model score and DMS score, followed by bootstrap confidence intervals over assays. ClinVar indels were reconstructed when protein-level sequence edits were unambiguous; unsupported frameshifts were excluded from the current scorer.

\subsection{Statistics}

ROC AUCs and mean Spearman correlations were reported with nonparametric bootstrap 95\% confidence intervals. For classifiers, macro-AUC averaged per-class AUCs across LOF, GOF and DN.

\section*{Data availability}

Public and license-gated resources are staged locally and on H200/OSS storage. A submission version should provide accession numbers, stable URLs and license notes.

\section*{Code availability}

R1 analysis code is under \texttt{Research1/}. A submission version should include a tagged release and frozen execution environment.

\section*{Acknowledgements}

TODO: add funding, compute and institutional acknowledgements.

\section*{Author contributions}

TODO: update after coauthor list is finalized.

\section*{Competing interests}

The authors declare no competing interests. TODO: confirm before submission.

\bibliography{references}

\end{document}
